\section{Posnania (wersja 1.1.1)}
\label{sec:posnania}

\subsection*{Okno i wolumen}
Okno eksperymentu: \textbf{2025-08-23 17{:}54{:}14–17{:}57{:}08} ($\approx$ \SI{2.9}{\minute}). 
Zarejestrowano \textbf{21} przepływy, wolumen $\approx$ \textbf{0.08 MiB} 
(\textbf{\emph{RX} $\approx$ 0.06 MiB}, \textbf{\emph{TX} $\approx$ 0.03 MiB}).

\paragraph{Przebieg badania:} 
Scenariusz obejmował około \textbf{15 minut interakcji} z aplikacją, w tym: 
\emph{nasłuch tła, start aplikacji, zmianę hasła, logowanie, dodanie rejestracji, nawigację po menu, wylogowanie, próbę błędnego logowania, ponowne logowanie, krótką pracę w trybie offline i online, ponowne wylogowanie oraz około 10 minut działania aplikacji w tle}.

\subsection{Wolumen i liczność wg protokołu}
\begin{table}[H]
  \centering
  \caption{Podsumowanie wg protokołu (Posnania, ta sesja)}
  \label{tab:posnania-proto}
  \begin{tabular}{lrrrr}
    \toprule
    \textbf{Protokół} & \textbf{Flows} & \textbf{MiB (total)} & \textbf{B wysłane} & \textbf{B odebrane} \\
    \midrule
    HTTPS & 9  & 0.08 & 27{,}627 & 57{,}363 \\
    DNS   & 12 & 0.00 & 260      & 468 \\
    \midrule
    \textbf{Razem} & \textbf{21} & \textbf{0.08} & \textbf{27{,}887} & \textbf{57{,}831} \\
    \bottomrule
  \end{tabular}
\end{table}

\paragraph{Obserwacja:} 
Ruch aplikacji jest w całości oparty o \textbf{\emph{HTTPS}} (brak \emph{QUIC} oraz \emph{HTTP} w badanym oknie czasowym).

\subsection{Największe cele (IP/port) wg wolumenu}
\begin{table}[H]
  \centering
  \caption{Największe kierunki ruchu (Posnania, ta sesja)}
  \label{tab:posnania-topdst}
  \begin{tabular}{lllrr}
    \toprule
    \textbf{IP docelowe} & \textbf{Port} & \textbf{Proto} & \textbf{Flows} & \textbf{MiB} \\
    \midrule
    104.21.42.187  & 443 & HTTPS & 1  & 0.04 \\
    34.120.195.249 & 443 & HTTPS & 5  & 0.02 \\
    52.26.132.221  & 443 & HTTPS & 2  & 0.01 \\
    172.67.208.92  & 443 & HTTPS & 1  & 0.01 \\
    10.215.173.2   & 53  & DNS   & 12 & 0.00 \\
    \bottomrule
  \end{tabular}
\end{table}

\paragraph{Obserwacja:}
Cele to wyłącznie port \textbf{443/\emph{TCP}} (\emph{HTTPS}) oraz standardowe \textbf{\emph{DNS}/53}. Brak nietypowych portów.

\subsection{Najczęstsze hosty (liczba przepływów)}
\begin{table}[H]
  \centering
  \caption{Top hosty (Posnania, liczba przepływów)}
  \label{tab:posnania-hosts}
  \begin{tabular}{lr}
    \toprule
    \textbf{Host} & \textbf{Flows} \\
    \midrule
    \texttt{api.yourb.app} & 10 \\
    \texttt{o4505784913690624.ingest.sentry.io} & 6 \\
    \texttt{m.stripe.com} & 3 \\
    \texttt{updates.yourb.app} & 2 \\
    \bottomrule
  \end{tabular}
\end{table}

\paragraph{Obserwacja:}
Widoczne są domeny \texttt{yourb.app} (\emph{API}/aktualizacje), \texttt{sentry.io} (raportowanie błędów) oraz \texttt{stripe.com} (płatności).

\subsection{Obserwacje jakościowe}
\begin{itemize}
  \item \textbf{Brak \emph{QUIC}/\emph{HTTP/3}.} Cała komunikacja po 443/\emph{TCP} (\emph{TLS}).
  \item \textbf{Brak \emph{HTTP}.} W badanej sesji nie obserwowano niezaszyfrowanych żądań \emph{HTTP}.
  \item \textbf{Integracje zewnętrzne.} Ruch do \texttt{sentry.io} (crash reporting) i \texttt{stripe.com} (płatności) jest zgodny z funkcjonalnością aplikacji.
\end{itemize}

\subsection{Wyniki z Exodus Privacy \cite{ExodusPrivacy}}
Raport dla \texttt{com.yourbuildingapp.app.posnania} (wersja 1.1.1, Google Play) wskazuje obecność \textbf{1 trackera} (\emph{Sentry} — crash reporting) oraz \textbf{18 uprawnień}. 

\textbf{Uprawnienia — przykłady kluczowe:}
\begin{itemize}
  \item \textbf{Kamera i multimedia:} \texttt{CAMERA}, \texttt{READ\_EXTERNAL\_STORAGE}, \texttt{WRITE\_EXTERNAL\_STORAGE}, \texttt{READ\_MEDIA\_IMAGES}, \texttt{READ\_MEDIA\_VIDEO}.
  \item \textbf{Tożsamość/telefonia:} \texttt{READ\_PHONE\_STATE}.
  \item \textbf{Interfejs/system:} \texttt{SYSTEM\_ALERT\_WINDOW}, \texttt{RECEIVE\_BOOT\_COMPLETED}, \texttt{SCHEDULE\_EXACT\_ALARM}, \texttt{POST\_NOTIFICATIONS}.
  \item \textbf{Sieć:} \texttt{INTERNET}, \texttt{ACCESS\_NETWORK\_STATE}, \texttt{ACCESS\_WIFI\_STATE}.
  \item \textbf{Biometria:} \texttt{USE\_BIOMETRIC}, \texttt{USE\_FINGERPRINT}.
\end{itemize}

\emph{Spostrzeżenia:} integracja z Sentry jest zgodna z profilem aplikacji, jednak szeroki zestaw uprawnień (kamera, biometria, alarmy systemowe) może być postrzegany jako nadmiarowy względem prostej funkcji aplikacji centrum handlowego.

\subsection{Subiektywne spostrzeżenia}
\paragraph{Obserwowalność treści.}
Ze względu na \emph{TLS} treści żądań/odpowiedzi nie są dostępne; brak oznak jawnego przesyłania danych.

\paragraph{Hipotezy dot. bezpieczeństwa (hipoteza wymagająca osobnych badań w celu klaryfikacji).}
Podczas korzystania z aplikacji spostrzeżono możliwość wycieku informacji danych wrażliwych. Dodawanie tablic rejestracyjnych pojazdów nie wymaga wykazania własności samochodu. Dodać można każdą rejestrację, a następnie śledzić zarejestrowane na nią wjazdy i wyjazdy z parkingu.
Na podstawie samego ruchu sieciowego \emph{nie można} potwierdzić „otwartej”/słabo szyfrowanej bazy danych — wymaga to analizy on-device (np. \texttt{SQLite/Room}, manifest, reverse engineering APK) oraz testów dynamicznych (np. \texttt{Frida}, \texttt{mitmproxy}, sprawdzenie pinu certyfikatu).

\subsection{Podsumowanie}
W badanej sesji Posnania wykazuje:
\begin{enumerate}
\item ruch wyłącznie po \emph{HTTPS} (bez \emph{QUIC}), niewielki wolumen ($\approx 0.08$ MiB) i zapytania \emph{DNS},
\item integrację z \texttt{sentry.io} (telemetria błędów) i \texttt{stripe.com} (płatności),
\item brak śladów HTTP w badanym oknie czasowym; dalsze badania wymagane, by ocenić bezpieczeństwo lokalnej bazy danych.
\end{enumerate}
